\section{Canonical pullback comparisons for the actual Cartier ratio sheaf}
\begin{lemma}[Canonical whole sheaf tensor and Cartier ratio pullback]
\label{lem:cartier-ratio-canonical-pullback}
\uses{lem:cartier-actual-invertible-module-sheaf,lem:actual-cartier-ratio-line-sheaf}
\lean{AlgebraicGeometry.Scheme.Modules.canonicalTensorPullbackIso,AlgebraicGeometry.Scheme.IdealSheafData.ratioPullbackIso}
\leanok
The pullback of the actual sheafified tensor along any genuine scheme
morphism is canonically the sheafified tensor of the actual pullbacks.
For a genuine Cartier ideal in the denominator, the actual ratio sheaf
pullback is canonically the tensor of the numerator pullback and the
true sheaf dual of the denominator pullback.
\end{lemma}
\begin{proof}
Use the existing genuine comparison between the localized coherent
monoidal tensor and the explicit sheafified tensor. Pull back its inverse,
then apply the inverse tensorator of the actual canonical monoidal
pullback functor, and finally the explicit tensor comparison on the
original source scheme. Every map is the actual proved sheafification
or monoidal pullback comparison; no arbitrary isomorphism is chosen.
For the ratio apply this construction to the original numerator module
and original actual sheaf dual of the denominator module. The Cartier
ideal's genuine module sheaf is invertible, so the existing canonical
dual-pullback map is an isomorphism. Tensor its actual isomorphism with
the numerator identity. This gives precisely the stated original ratio
pullback isomorphism, for arbitrary original scheme morphisms.
\end{proof}

\begin{lemma}[Actual canonical ratio comparison under whole Cartier invariance]
\label{lem:cartier-ratio-whole-invariant-comparison}
\uses{lem:cartier-ratio-canonical-pullback,lem:cartier-invariant-module-iso}
\lean{AlgebraicGeometry.Scheme.IdealSheafData.IsInvertible.invariantRatioModuleIso}
\leanok
If two actual flat scheme morphisms preserve the two entire Cartier
ideal comaps, their genuine ratio module pullbacks have an actual
canonical isomorphism.
\end{lemma}
\begin{proof}
Use the preceding actual ratio pullback comparison for the first
morphism. Tensor the genuine canonical numerator comparison with the
true contravariant sheaf dual of the inverse denominator comparison.
Then compose with the inverse actual ratio pullback comparison for the
second morphism. The denominator inverse is required because the actual
sheaf dual is contravariant. Every factor is an actual original sheaf
isomorphism constructed from the genuine full geometric ideal equality.
No module comparison input is used. Group unit and cocycle coherence of
this tensor-dual ratio, its model metric/polarization and nefness are
separate subsequent assertions.
\end{proof}


\begin{lemma}[Actual whole invariant ratio identity and transitivity]
\label{lem:cartier-ratio-whole-invariant-transitivity}
\uses{lem:cartier-ratio-whole-invariant-comparison,lem:cartier-actual-tensor-dual-functorial,lem:cartier-invariant-module-coherence}
\lean{AlgebraicGeometry.Scheme.IdealSheafData.IsInvertible.invariantRatioModuleIso_refl,AlgebraicGeometry.Scheme.IdealSheafData.IsInvertible.invariantRatioModuleIso_trans}
\leanok
The actual invariant Cartier ratio module comparison has identity and
transitivity laws on the whole original source.
\end{lemma}
\begin{proof}
For identity, both actual factor comparisons are the original identity.
The genuine sheaf dual and explicit sheafified tensor preserve those
identities, so the two original ratio pullback comparisons cancel.
For transitivity, the inverse middle ratio pullback comparison cancels
with the identical original forward comparison. True tensor
functoriality composes the numerator comparisons and the two reversed
dual denominator comparisons. Genuine contravariant dual functoriality
identifies the latter with the dual of the inverse composite denominator
comparison. The already proved whole factor transitivity laws now give
exactly the original ratio comparison between the first and last maps.
All maps are the actual constructed isomorphisms; no ratio
identity or transitivity premise is used.
\end{proof}
